% RQ3 Methodology: Pipeline Experiments

\subsection{Experimental Design}

We conduct a factorial experiment varying three pipeline parameters:

\begin{table}[htbp]
\centering
\caption{RQ3 Experiment configuration (Experiment 05)}
\label{tab:rq3-experiment}
\begin{tabular}{@{}llll@{}}
\toprule
Factor & Levels & Values & Description \\
\midrule
Temp-check & 3 & None, 20\%, 40\% & Threshold to advance \\
Fast-track & 3 & None, 70\%, 85\% & Support threshold for fast-track \\
Expiry & 3 & 30, 60, 90 days & Maximum proposal lifetime \\
\bottomrule
\end{tabular}
\end{table}

Total configurations: $3 \times 3 = 9$ (temp-check $\times$ fast-track, with expiry as secondary)

Each configuration runs 30 times with different seeds.

Total runs: 270

\subsection{Parameter Selection}

\subsubsection{Temp-Check Levels}

\begin{itemize}
    \item \textbf{None}: All proposals proceed directly to formal voting
    \item \textbf{20\% threshold}: Low bar; filters only proposals with minimal support
    \item \textbf{40\% threshold}: Higher bar; filters controversial or low-awareness proposals
\end{itemize}

\subsubsection{Fast-Track Levels}

\begin{itemize}
    \item \textbf{None}: All proposals follow standard timeline
    \item \textbf{70\% threshold}: Proposals with supermajority early support fast-tracked
    \item \textbf{85\% threshold}: Only near-unanimous proposals fast-tracked
\end{itemize}

\subsubsection{Expiry Levels}

\begin{itemize}
    \item \textbf{30 days}: Tight window; proposals must resolve quickly
    \item \textbf{60 days}: Moderate window; allows extended deliberation
    \item \textbf{90 days}: Loose window; accommodates complex multi-round processes
\end{itemize}

\subsection{Baseline Configuration}

Common baseline parameters:

\begin{itemize}
    \item Members: 200 agents
    \item Proposal frequency: 0.5/day
    \item Voting period: 7 days
    \item Quorum: 4\%
    \item Simulation length: 2,000 steps
    \item Temp-check turnout: 30\% of voting population
    \item Temp-check to formal vote correlation: $\rho = 0.7$
\end{itemize}

\subsection{Hypotheses}

\begin{itemize}
    \item \textbf{H3.1}: Higher temp-check thresholds reduce formal voting load but increase false negative rate (good proposals filtered)
    \item \textbf{H3.2}: Fast-tracks reduce mean time-to-decision for qualifying proposals
    \item \textbf{H3.3}: Tighter expiry windows increase abandonment rate for complex proposals
    \item \textbf{H3.4}: There exists an optimal temp-check threshold balancing throughput and quality
\end{itemize}

\subsection{Metrics}

Primary metrics for RQ3:

\begin{description}
    \item[Throughput] Proposals resolved per time unit
    \item[Pass Rate] Fraction of proposals that pass
    \item[Time-to-Decision] Mean time from creation to resolution
    \item[Abandonment Rate] Proposals abandoned before resolution
    \item[Expiry Rate] Proposals expiring without resolution
    \item[Quorum Rate] Fraction of formal votes achieving quorum
\end{description}

Secondary metrics:

\begin{description}
    \item[Temp-check conversion] Fraction passing temp-check
    \item[Fast-track rate] Fraction of proposals fast-tracked
    \item[Stage durations] Time spent at each pipeline stage
\end{description}

\subsection{Statistical Analysis}

\subsubsection{Main Effects}

ANOVA for each factor's effect on outcome metrics.

\subsubsection{Interaction Effects}

Two-way interactions between temp-check and fast-track settings.

\subsubsection{Optimal Configuration}

Multi-objective optimization to identify configurations on Pareto frontier of:
\begin{itemize}
    \item Maximize throughput
    \item Maximize pass rate
    \item Minimize abandonment
    \item Minimize time-to-decision
\end{itemize}
